\chapter{Quadratic Residues and Reciprocity}\label{ch:reciprocity}

\begin{center}\itshape
Is $219$ a square modulo the prime $383$? A direct search might square $191$ numbers.\\
With it, the answer (yes) takes four lines. Gauss called it the \emph{golden theorem} and returned to it again and again.
\end{center}

Euler noticed it, Legendre half-proved it, and on 8 April 1796, a few weeks before his nineteenth birthday, Gauss proved it: whether a prime $p$ is a
square modulo another prime $q$ is tied, by a simple rule, to whether $q$ is a square modulo $p$. There is no obvious reason why the two questions should
have anything to do with each other, and that is exactly why the theorem is so loved. Gauss called it the \emph{golden theorem}, and returned to it all his
life; many different proofs are now known. Here we give Eisenstein's, which counts lattice points in a rectangle, and an exercise gives a second one with Gauss sums.

\portrait{gauss}{Carl Friedrich Gauss (1777--1855). Portrait by Christian Albrecht Jensen, 1840. Public domain, via Wikimedia Commons.}

\begin{reminder}[Three questions, in the right order]
First: is a non-zero class a square? Euler's criterion answers by a modular
power. Second: how can we decide faster, or see the pattern as the prime
changes? Reciprocity replaces one prime by a smaller remainder. Third: if a
square root exists, how do we find it? That is a separate problem; a Legendre
symbol answers yes or no and does not hand over the root. We will keep these
questions apart, then connect them in a worked example.
\end{reminder}

\section{Quadratic residues and Euler's criterion}

Let $p$ be an odd prime. An integer $a$ with $p\nmid a$ is a \emph{quadratic residue} modulo $p$ if $x^2\equiv a\pmod p$ has a solution, and a \emph{non-residue} otherwise.

\begin{definition}
The \emph{Legendre symbol} is $\left(\frac ap\right)=1$ if $a$ is a quadratic residue, $-1$ if it is a non-residue, and $0$ if $p\mid a$.
\end{definition}

\begin{think}
List the squares modulo $11$: how many distinct non-zero classes do you get? Then compute $a^5\bmod11$ for $a=1,\dots,10$. Which values occur, and which $a$ give which value? Make a conjecture before reading the theorem.
\end{think}

\begin{theorem}[Euler's criterion]\label{thm:eulercrit}
$\left(\dfrac ap\right)\equiv a^{(p-1)/2}\pmod p$. Consequently the symbol is completely multiplicative in $a$, there are exactly $\frac{p-1}2$ residues and $\frac{p-1}2$ non-residues, and
\[
\left(\frac{-1}p\right)=(-1)^{(p-1)/2}:
\]
$-1$ is a square modulo $p$ if and only if $p\equiv1\pmod4$.
\end{theorem}
\begin{proof}
For $p\nmid a$, let $g$ be a primitive root (Chapter~\ref{ch:orders}) and write $a\equiv g^i\pmod p$. Then $a$ is a square if and only if $i$ is even (if $a=g^{2j}$, $x=g^j$ works; if $x=g^j$ then $x^2=g^{2j}$ and $2j\equiv i\pmod{p-1}$ forces $i$ even). Also $a^{(p-1)/2}=g^{i(p-1)/2}$, which is $1$ if $i$ is even and $g^{(p-1)/2}=-1$ if $i$ is odd (it squares to $1$ and is not $1$). Half of the exponents $0,\dots,p-2$ are even.
For $p\mid a$, both sides of the criterion are $0$, since $(p-1)/2\ge1$.
For any $a,b$, applying the criterion to $ab$ gives
$\left(\frac{ab}p\right)\equiv
\left(\frac ap\right)\left(\frac bp\right)\pmod p$.
If neither factor is zero these quantities belong to $\{1,-1\}$, whose
classes modulo the odd prime $p$ are distinct; if a factor is zero both
sides vanish. Thus the multiplicativity is an equality of integers, not just
a congruence. Finally apply the criterion with $a=-1$ to get the supplementary law.
\end{proof}

\paragraph{Why exactly two roots?}
If $x^2\equiv y^2\pmod p$, then $p\mid(x-y)(x+y)$, so $x\equiv y$ or
$x\equiv-y$. Thus each non-zero square has exactly two roots, a pair
$\{x,-x\}$. There is no overlap inside such a pair because $p$ is odd and
$x\ne0$. The $p-1$ units therefore produce $(p-1)/2$ squares. This also
explains why zero is excluded from the word ``residue'' here: it has only one
root and the Legendre symbol treats it separately.

Modulo $11$, squaring $1,2,3,4,5$ gives $1,4,9,5,3$. The five non-zero
residues are therefore $1,3,4,5,9$. A compact experiment already displays the
criterion: $3^5\equiv1$, but $2^5=32\equiv-1$ modulo $11$.
Multiplying two non-residues gives a residue, since $(-1)(-1)=1$ at the level
of symbols. Addition has no such rule: $1$ and $3$ are residues but their sum
$4$ is too, while $1+1=2$ is a non-residue. The symbol is a multiplicative
sign, not an arithmetic fortune-teller.

\subsection*{A first character, and its cancellation}
Write $\chi(a)=\left(\frac ap\right)$. The identities just proved say that
$\chi$ is periodic modulo $p$, multiplicative, and zero precisely on multiples
of $p$. Such a function is called a \emph{Dirichlet character modulo $p$}.
Here we need only this particular real-valued example, not a general theory.

Its sum over a full residue system is zero. One proof counts the equally
numerous signs. Another exposes a useful mechanism: choose a non-residue $b$.
Multiplication by $b$ permutes the residues, so
\[
S:=\sum_{a=0}^{p-1}\chi(a)
=\sum_{a=0}^{p-1}\chi(ba)=\chi(b)S=-S,
\]
whence $S=0$. The cancellation is exact although the individual signs may
look irregular. It will let us calculate an expectation in
Exercise~\ref{ex:t1-residue-chance} without assuming independent signs.

\section{Gauss' lemma}

\begin{theorem}[Gauss' lemma]\label{thm:gausslemma}
Let $p\nmid a$, and for $1\le k\le\frac{p-1}2$ let $r_k$ be the residue of $ka$ in $\left]-\frac p2,\frac p2\right[$. If $\mu$ is the number of negative $r_k$, then $\left(\frac ap\right)=(-1)^\mu$.
\end{theorem}
\begin{proof}
The $|r_k|$ lie in $\{1,\dots,\frac{p-1}2\}$ and are distinct: $ka\equiv\pm la$ implies $k\equiv\pm l$, so $k=l$ in this range. So they are a permutation of $1,\dots,\frac{p-1}2$. Multiplying all the $ka$:
$a^{(p-1)/2}\left(\frac{p-1}2\right)!\equiv\prod_kr_k=(-1)^\mu\left(\frac{p-1}2\right)!\pmod p$. Cancel the factorial and apply Euler's criterion.
\end{proof}

\paragraph{A signed permutation, seen explicitly.}
Take $p=11$ and $a=3$. Reduce $3,6,9,12,15$ to the symmetric interval
$(-11/2,11/2)$; the representatives are
\[
 3,-5,-2,1,4.
\]
Their absolute values are a permutation of $1,2,3,4,5$, and there are two
minus signs. Hence $\left(\frac3{11}\right)=1$, agreeing with
$5^2\equiv3\pmod{11}$. For $a=2$, the representatives
$2,4,-5,-3,-1$ have three minus signs, so
$\left(\frac2{11}\right)=-1$.

The key cancellation in the proof is legal because none of
$1,\dots,(p-1)/2$ is divisible by $p$; the factorial is a unit.
To justify the distinctness carefully, $k\equiv-l\pmod p$ would require
$k+l=p$, impossible since $2\le k+l\le p-1$.
Gauss has turned a question about a square root into the parity of a count.
We count signs, not the size of the representatives.

\begin{corollary}\label{cor:two}
$\left(\dfrac2p\right)=(-1)^{(p^2-1)/8}$: $2$ is a square modulo $p$ if and only if $p\equiv\pm1\pmod8$.
\end{corollary}
\begin{proof}
For $a=2$, the numbers $2k$ with $1\le k\le\frac{p-1}2$ lie in $[2,p-1]$, and $r_k<0$ exactly when $2k>\frac p2$. So $\mu=\frac{p-1}2-\left\lfloor\frac p4\right\rfloor$. Checking $p=8m+1,3,5,7$ gives $\mu=2m,2m+1,2m+1,2m+2$, whose parity is that of $\frac{p^2-1}8$.
\end{proof}

\section{The law of quadratic reciprocity}

\begin{think}
Make a table: for the primes $p,q\in\{3,5,7,11,13\}$, $p\ne q$, record whether $p$ is a square modulo $q$. Now compare the entry for $(p,q)$ with the entry for $(q,p)$. When are they equal, and when do they differ? Look at $p$ and $q$ modulo $4$.
\end{think}

\begin{theorem}[quadratic reciprocity; Gauss, 1796]\label{thm:qr}
For distinct odd primes $p$ and $q$,
\[
\left(\frac pq\right)\left(\frac qp\right)=(-1)^{\frac{p-1}2\cdot\frac{q-1}2}.
\]
That is, $\left(\frac pq\right)=\left(\frac qp\right)$ unless $p\equiv q\equiv3\pmod4$, in which case $\left(\frac pq\right)=-\left(\frac qp\right)$.
\end{theorem}

\begin{figure}[ht]
\centering
\fig{ch06-reciprocity-table}
\caption{Each cell records $\left(\frac pq\right)$: row prime $p$, column
prime $q$. A plus means a non-zero square, a minus a non-square, and $0$
on the diagonal records divisibility. Compare cells reflected across the
diagonal: they disagree exactly when both primes are $3$ modulo $4$.
Signs carry the information even in black-and-white print.}
\label{fig:reciprocitytable}
\end{figure}

We give Eisenstein's proof, which counts lattice points. A second proof, with Gauss sums, is Exercise~\ref{ex:gausssum}.

\begin{lemma}[Eisenstein]\label{lem:eisenstein}
For odd $a$ with $p\nmid a$, $\left(\dfrac ap\right)=(-1)^{T(a,p)}$ where $T(a,p)=\displaystyle\sum_{k=1}^{(p-1)/2}\left\lfloor\frac{ka}p\right\rfloor$.
\end{lemma}
\begin{proof}
Write $ka=p\left\lfloor\frac{ka}p\right\rfloor+s_k$ with $0<s_k<p$. Let $A$ be the sum of the $s_k<\frac p2$ and $B$ the sum of the $\mu$ values $s_k>\frac p2$. In the proof of Gauss' lemma, the numbers $s_k$ ($s_k<\frac p2$) and $p-s_k$ ($s_k>\frac p2$) are $1,\dots,\frac{p-1}2$ in some order. Let $S=\sum_{k=1}^{(p-1)/2}k=\frac{p^2-1}8$. Then
\[
aS=p\,T(a,p)+A+B,\qquad S=A+\mu p-B.
\]
Subtracting, $(a-1)S=p\left(T(a,p)-\mu\right)+2B$. Modulo $2$, with $a$ and $p$ odd: $0\equiv T(a,p)-\mu$. Conclude with Gauss' lemma.
\end{proof}

\begin{figure}[ht]
\centering
\fig{ch06-eisenstein}
\caption{Eisenstein's count for $p=7$ and $q=11$. The $3\times5=15$ lattice points of the rectangle are split by the diagonal $y=\frac{11}7x$, which passes through none of them, into $T(11,7)=8$ and $T(7,11)=7$. So $\left(\frac{11}7\right)\left(\frac7{11}\right)=(-1)^{15}=-1$: indeed $11\equiv2^2$ is a square modulo $7$, but $7$ is not a square modulo $11$.}
\label{fig:eisenstein}
\end{figure}

\begin{think}
Before reading the proof, look at Figure~\ref{fig:eisenstein}. Count the points below the diagonal column by column, and those above it row by row. What do you get, and why does Eisenstein's lemma turn this picture into reciprocity?
\end{think}

\begin{proof}[Proof of Theorem~\ref{thm:qr}]
Count the points $(x,y)\in\Z^2$ with $1\le x\le\frac{p-1}2$ and $1\le y\le\frac{q-1}2$: there are $\frac{p-1}2\cdot\frac{q-1}2$. None lies on the line $qx=py$ (it would force $p\mid x$). Below the line ($y<\frac{qx}p$; note $\frac{qx}p<\frac q2$, so the bound $y\le\frac{q-1}2$ is automatic there) there are $\sum_x\left\lfloor\frac{qx}p\right\rfloor=T(q,p)$ points, and above it, symmetrically, $T(p,q)$. So $T(q,p)+T(p,q)=\frac{p-1}2\cdot\frac{q-1}2$, and Eisenstein's lemma gives the theorem.
\end{proof}

\paragraph{Why Eisenstein assumes that $a$ is odd.}
His proof ends by discarding $(a-1)S$ modulo $2$; this is even when $a$ is
odd. That assumption cannot simply be dropped. For $a=2,p=3$, the floor sum
is $\lfloor2/3\rfloor=0$, predicting a plus sign, whereas $2$ is not a square
modulo $3$. We handle factors of $2$ with the supplementary law, then apply
reciprocity only to odd numbers. Every hypothesis is doing a job.

\begin{example}
$383$ is prime and $219=3\cdot73$. Since $383\equiv3\pmod4$ and $3\equiv3\pmod4$: $\left(\frac3{383}\right)=-\left(\frac{383}3\right)=-\left(\frac23\right)=1$. Since $73\equiv1\pmod4$: $\left(\frac{73}{383}\right)=\left(\frac{383}{73}\right)=\left(\frac{18}{73}\right)=\left(\frac2{73}\right)\left(\frac9{73}\right)=1$ as $73\equiv1\pmod8$. So $\left(\frac{219}{383}\right)=1$.
\end{example}

\paragraph{A complete calculation, including the root.}
The previous example gives an existence decision. As a short alternative,
$18=2\cdot3^2$ modulo $73$, and $73\equiv1\pmod8$, so the only non-trivial
signs are already resolved by small remainders. What about producing a square
root of $219$ modulo $383$?

For a prime $p\equiv3\pmod4$ and a non-zero residue $a$, set
$r=a^{(p+1)/4}$. Then Euler's criterion gives
\[
 r^2=a^{(p+1)/2}=a\,a^{(p-1)/2}\equiv a\pmod p.
\]
Here $(383+1)/4=96$. Repeated squaring gives
\[
\begin{gathered}
219^2\equiv86,\quad219^4\equiv119,\quad219^8\equiv373,\\
219^{16}\equiv100,\quad219^{32}\equiv42,\quad219^{64}\equiv232
\pmod{383}.
\end{gathered}
\]
Thus $219^{96}\equiv232\cdot42\equiv169$ and the two roots are $169$ and
$383-169=214$. A final check is pleasantly short:
$169^2=28561=383\cdot74+219$. Existence, construction, verification:
each step has its own proof. The exponent formula is specific to
$p\equiv3\pmod4$; other primes require different methods.

\section{The Jacobi symbol}

For a positive odd integer $n=p_1\cdots p_r$ (primes with repetition), the \emph{Jacobi symbol} is $\left(\frac an\right)=\prod_i\left(\frac a{p_i}\right)$. It is multiplicative in $a$ (as each Legendre symbol is) and in $n$ (by definition). For $n=1$ the product is empty and the symbol is $1$, for every $a$.

\begin{proposition}\label{prop:jacobi}
For odd coprime $m,n>0$:
\[
\begin{gathered}\left(\frac{-1}n\right)=(-1)^{\frac{n-1}2},\qquad\left(\frac2n\right)=(-1)^{\frac{n^2-1}8},\qq\left(\frac mn\right)\left(\frac nm\right)=(-1)^{\frac{m-1}2\cdot\frac{n-1}2}.\end{gathered}
\]
\end{proposition}
\begin{proof}
For odd $a,b$: $\frac{ab-1}2-\frac{a-1}2-\frac{b-1}2=\frac{(a-1)(b-1)}2$ is even, and $\frac{a^2b^2-1}8-\frac{a^2-1}8-\frac{b^2-1}8=\frac{(a^2-1)(b^2-1)}8$ is even (each factor is divisible by $8$). So $u(n)=\frac{n-1}2$ and $v(n)=\frac{n^2-1}8$ are additive modulo $2$ over products of odd numbers. The first two formulas hold for primes (Theorem~\ref{thm:eulercrit}, Corollary~\ref{cor:two}), and both sides are multiplicative in $n$. For the third, with $m=\prod_iq_i$ and $n=\prod_jp_j$ (all distinct primes across the two lists, since $\gcd(m,n)=1$),
$\left(\frac mn\right)\left(\frac nm\right)=\prod_{i,j}\left(\frac{q_i}{p_j}\right)\left(\frac{p_j}{q_i}\right)=(-1)^{\sum_{i,j}u(q_i)u(p_j)}=(-1)^{u(m)u(n)}$, using Theorem~\ref{thm:qr} and additivity of $u$.
\end{proof}

This gives a Euclid-like algorithm computing $\left(\frac an\right)$ \emph{without factoring}. Beware: $\left(\frac an\right)=1$ does not imply that $a$ is a square modulo a composite $n$ (e.g.\ $\left(\frac2{15}\right)=1$, but $2$ is not a square modulo $3$).

\paragraph{Computing a Jacobi symbol without factoring.}
Here is the algorithm behind the claim. Start with a positive odd denominator
$n$ and a sign $s=1$. At each stage reduce $a$ modulo $n$. If $a=0$, stop:
the answer is $0$ when $n>1$, and $s$ when $n=1$. If $a=1$, stop with
answer $s$. Otherwise remove every factor of $2$ from $a$, multiplying $s$
each time by $(-1)^{(n^2-1)/8}$. If this leaves $a=1$, stop with answer
$s$. The remaining $a>1$ is odd: swap $a$ and $n$, flipping $s$ when both
are $3$ modulo $4$, then repeat. After the reduction, the new denominator
is the old $a<n$, so the process terminates. The rule at $n=1$ is the
empty-product convention. The same reductions are Euclid's algorithm with
a few signs attached.
For non-coprime inputs, reciprocity still gives the required equality with
both symbols equal to $0$: an odd prime divides both numerator and denominator.
Removing factors of $2$ and taking Euclidean remainders preserve their gcd,
so this case ends at the zero rule with denominator greater than $1$.
This explains why the algorithm does not need a separate coprimality test.

For example, every swap in the following chain has positive sign because
at least one of its two odd numbers is $1$ modulo $4$:
\[
\begin{aligned}
\left(\frac{1001}{9907}\right)
 &=\left(\frac{898}{1001}\right)
 =\left(\frac{449}{1001}\right)\\
 &=\left(\frac{103}{449}\right)
 =\left(\frac{37}{103}\right)
 =\left(\frac{29}{37}\right)
 =\left(\frac8{29}\right)=-1.
\end{aligned}
\]
Here $\left(\frac2{1001}\right)=1$ because $1001\equiv1\pmod8$,
and $\left(\frac2{29}\right)=-1$ because $29\equiv5\pmod8$.
No factorisation of $9907$ was needed. A negative answer rules out a square;
a positive answer for a composite denominator need not certify one. The two
non-residues $2$ modulo $3$ and $2$ modulo $5$ give
$\left(\frac2{15}\right)=(-1)(-1)=1$, while $x^2\equiv2\pmod{15}$ has
no solution. A Jacobi symbol has forgotten some local information.

\section{Exercises}

\begin{exercise}[$\star$]\label{ex:minus3}
Show that for a prime $p>3$, $-3$ is a square modulo $p$ if and only if $p\equiv1\pmod3$.
\end{exercise}

\begin{exercise}[$\star$]\label{ex:onemod4}
Show that every odd prime factor of $n^2+1$ is $\equiv1\pmod4$, and deduce that there are infinitely many primes $p\equiv1\pmod4$.
\end{exercise}

\begin{exercise}[$\star\star$]\label{ex:minus2}
For an odd prime $p$, show that $\left(\frac{-2}p\right)=1$ if and only if $p\equiv1$ or $3\pmod8$. Deduce that a prime of the form $x^2+2y^2$ is $2$ or $\equiv1,3\pmod8$.
\end{exercise}

\begin{exercise}[$\star\star$]\label{ex:mersfactor}
Let $p\equiv3\pmod4$ be prime with $q=2p+1$ also prime. Show that $q\mid2^p-1$. Deduce that $2^{11}-1$, $2^{23}-1$ and $2^{251}-1$ are composite.
\end{exercise}

\begin{exercise}[$\star\star$]\label{ex:consecutive}
Let $p$ be an odd prime. Show that $\displaystyle\sum_{a=1}^{p-1}\left(\frac{a(a+1)}p\right)=-1$. Deduce the number of $a\in\{1,\dots,p-2\}$ such that $a$ and $a+1$ are both quadratic residues.
\end{exercise}

\begin{exercise}[$\star\star$]\label{ex:pepin}
\emph{P\'epin's test.} Let $\mathcal F_n=2^{2^n}+1$ with $n\ge1$. Show that $\mathcal F_n$ is prime if and only if $3^{(\mathcal F_n-1)/2}\equiv-1\pmod{\mathcal F_n}$.
\end{exercise}

\begin{exercise}[$\star\star\star$]\label{ex:gausssum}
\emph{Reciprocity through Gauss sums.} Let $\zeta=e^{2i\pi/p}$, $p^*=\left(\frac{-1}p\right)p$, and $g=\sum_{a=1}^{p-1}\left(\frac ap\right)\zeta^a$.
\begin{enumerate}[label=(\alph*)]
\item Show that $g^2=p^*$. \emph{Hint: substitute $b=ac$ in $g^2=\sum_{a,b}\left(\frac{ab}p\right)\zeta^{a+b}$.}
\item Let $q\ne p$ be an odd prime. In the ring $\Z[\zeta]$, show that $g^q\equiv\left(\frac qp\right)g\pmod{q\Z[\zeta]}$.
\item Show that $q\Z[\zeta]\cap\Z=q\Z$, using that $\Phi_p$ is irreducible over $\Q$ (Volume~0, the polynomial exercise on prime cyclotomic irreducibility). Deduce $\left(\frac{p^*}q\right)=\left(\frac qp\right)$, and then Theorem~\ref{thm:qr}.
\item Show that $g=\sum_{x=0}^{p-1}\zeta^{x^2}$, the sum $G$ of Volume~0, the complex-number exercise on quadratic Gauss sums. What does (a) say about $G$?
\end{enumerate}
\end{exercise}

\begin{exercise}[$\star\star$]\label{ex:t1-residue-chance}
\emph{Residue signs and chance.} Choose $U$ uniformly from
$\{0,\dots,p-1\}$ for an odd prime $p$, and put
$X=\left(\frac Up\right)$, $Y=\left(\frac{U+1}p\right)$.
Find the distributions, expectations and variances of $X,Y$.
Using Exercise~\ref{ex:consecutive}, prove
$\operatorname{Cov}(X,Y)=-1/p$. Are the two signs independent?
\end{exercise}

\begin{beauty}
Quadratic reciprocity links two questions that have no visible connection: whether $p$ is a square modulo $q$, and whether $q$ is a square modulo $p$. Many proofs are known, each exposing a different structure. Its generalisations (cubic and biquadratic reciprocity, which live in the rings of Chapter~\ref{ch:gaussian}, Artin reciprocity, and the Langlands program) have driven number theory for two centuries. It also powers the Solovay--Strassen primality test and the computation of square roots modulo $p$.
\end{beauty}

\section{Hints}

\begin{hint}{ex:minus3}
$\left(\frac{-3}p\right)=\left(\frac{-1}p\right)\left(\frac3p\right)$. Apply Euler's criterion to the first factor and reciprocity to the second: the signs cancel.
\end{hint}
\begin{hint}{ex:onemod4}
$p\mid n^2+1$ means that $-1$ is a square modulo $p$. For infinitude, take $n=2p_1\cdots p_k$.
\end{hint}
\begin{hint}{ex:minus2}
Multiply the two supplementary laws and tabulate $p\bmod8$. For the second part, if $p=x^2+2y^2$ then $p\nmid y$, and $(xy^{-1})^2\equiv-2$.
\end{hint}
\begin{hint}{ex:mersfactor}
$q\mid2^p-1$ means $2^{(q-1)/2}\equiv1\pmod q$, that is, $2$ is a square modulo $q$. What is $q\bmod8$?
\end{hint}
\begin{hint}{ex:consecutive}
For $a\ne0$, $\left(\frac{a(a+1)}p\right)=\left(\frac{1+a^{-1}}p\right)$. Which values does $1+a^{-1}$ take as $a$ runs over the units? For the count, express the indicator of ``$a$ and $a+1$ are residues'' with $\frac{1+\chi(a)}2\cdot\frac{1+\chi(a+1)}2$.
\end{hint}
\begin{hint}{ex:pepin}
If $\mathcal F_n$ is prime, compute $\left(\frac3{\mathcal F_n}\right)$ by reciprocity, using $\mathcal F_n\bmod3$ and $\mathcal F_n\bmod4$. For the converse, the congruence gives the order of $3$ modulo any prime factor $p$, and it must divide $p-1$.
\end{hint}
\begin{hint}{ex:gausssum}
(b) In characteristic $q$, $(x+y)^q\equiv x^q+y^q$, so $g^q\equiv\sum_a\left(\frac ap\right)\zeta^{aq}$; substitute $b=aq$. (c) Multiply (b) by $g$ and use (a); then both sides are integers. Do not cancel $g$ inside the quotient ring before proving that it is invertible there. Instead descend to an integer congruence and cancel $p^*$ modulo $q$. (d) Count the solutions of $x^2\equiv a$ with $1+\left(\frac ap\right)$.
\end{hint}

\begin{hint}{ex:t1-residue-chance}
Zero occurs once and each nonzero sign occurs $(p-1)/2$ times.
Sum $\chi(a)\chi(a+1)$ over all $a$; the missing term at $a=0$ is zero.
\end{hint}

\section{Solutions}

\begin{solution}{ex:minus3}
\idea Split $-3=(-1)\cdot3$ by multiplicativity: $\left(\frac{-1}p\right)$ is given by Euler's criterion and $\left(\frac3p\right)$ by reciprocity. The two signs $(-1)^{(p-1)/2}$ cancel, leaving $\left(\frac p3\right)$, which only depends on $p\bmod3$.
\writeup $\left(\frac{-3}p\right)=\left(\frac{-1}p\right)\left(\frac3p\right)=(-1)^{\frac{p-1}2}\cdot(-1)^{\frac{p-1}2}\left(\frac p3\right)=\left(\frac p3\right)$, by reciprocity with the exponent $\frac{p-1}2\cdot\frac{3-1}2=\frac{p-1}2$. And $\left(\frac p3\right)=1$ iff $p\equiv1\pmod3$ (the squares modulo $3$ are $0,1$).
\end{solution}

\begin{solution}{ex:onemod4}
\idea $p\mid n^2+1$ says exactly that $-1$ is a square modulo $p$, which Euler's criterion translates into $p\equiv1\pmod4$. For infinitude, imitate Euclid with $n^2+1$ in place of $n+1$, choosing $n$ divisible by all the known primes $\equiv1\pmod4$ (and by $2$, to make $n^2+1$ odd).
\writeup If $p\mid n^2+1$ with $p$ odd, $-1\equiv n^2$ is a square modulo $p$, so $p\equiv1\pmod4$ by Euler's criterion. If $p_1,\dots,p_k$ were all the primes $\equiv1\pmod4$, then $N=(2p_1\cdots p_k)^2+1>1$ is odd, and its prime factors are $\equiv1\pmod4$ but differ from all the $p_i$ (which would divide $1$).
\end{solution}

\begin{solution}{ex:minus2}
\idea Multiply the two supplementary laws, $\left(\frac{-1}p\right)$ and $\left(\frac2p\right)$, and tabulate modulo $8$. For the second part: $p=x^2+2y^2$ gives $x^2\equiv-2y^2$, so $-2$ is a square as soon as $y$ is invertible.
\writeup $\left(\frac{-2}p\right)=(-1)^{\frac{p-1}2}(-1)^{\frac{p^2-1}8}$. For $p\equiv1,3,5,7\pmod8$ the two signs are $(+,+),(-,-),(+,-),(-,+)$, so the product is $1$ exactly for $p\equiv1,3$. If $p=x^2+2y^2$ is odd, then $p\nmid y$ (else $p\mid x$ and $p^2\mid p$), and $\left(xy^{-1}\right)^2\equiv-2$.
\end{solution}

\begin{solution}{ex:mersfactor}
\idea $q\mid2^p-1$ means $2^{(q-1)/2}\equiv1\pmod q$, which by Euler's criterion means $2$ is a square modulo $q$. And $\left(\frac2q\right)$ only depends on $q\bmod8$: here $q=2p+1$ with $p\equiv3\pmod4$ gives $q\equiv7\pmod8$.
\writeup $q=2p+1\equiv7\pmod8$, so $\left(\frac2q\right)=1$ (Corollary~\ref{cor:two}) and Euler's criterion gives $2^{(q-1)/2}=2^p\equiv1\pmod q$. For $p=11,23,251$ the numbers $q=23,47,503$ are prime, and $q<2^p-1$, so $q$ is a proper factor.
\end{solution}

\begin{solution}{ex:consecutive}
\idea Sums of Legendre symbols of polynomials become tractable after a substitution that turns the polynomial into something linear. Factor out the square $a^2$: $\chi(a(a+1))=\chi(1+a^{-1})$, and $1+a^{-1}$ runs over almost all residues. For the count, the indicator of ``$a$ is a residue'' is $\frac{1+\chi(a)}2$, so the count of pairs is a sum of products of such indicators, which expands into sums we know.
\writeup For $1\le a\le p-1$, $\left(\frac{a(a+1)}p\right)=\left(\frac{a^2(1+a^{-1})}p\right)=\left(\frac{1+a^{-1}}p\right)$. As $a$ runs over the non-zero residues, $1+a^{-1}$ runs over all residues except $1$. Since $\sum_{c=0}^{p-1}\left(\frac cp\right)=0$ (equally many residues and non-residues), the sum is $0-\left(\frac1p\right)=-1$.
Let $N$ be the number of $a\in[1,p-2]$ with both $a$, $a+1$ residues. With $\chi=\left(\frac\cdot p\right)$, $\frac14(1+\chi(a))(1+\chi(a+1))$ is $1$ for such $a$ and $0$ for other $a\in[1,p-2]$. Summing over $a\in[1,p-2]$: $\sum\chi(a)=-\chi(p-1)=-\chi(-1)$, $\sum\chi(a+1)=-\chi(1)=-1$, and $\sum\chi(a)\chi(a+1)=-1-\chi(p-1)\chi(p)=-1$. So $N=\frac14\left(p-2-\chi(-1)-1-1\right)=\frac{p-4-\left(\frac{-1}p\right)}4$: $\frac{p-5}4$ if $p\equiv1\pmod4$ and $\frac{p-3}4$ if $p\equiv3\pmod4$.
\end{solution}

\begin{solution}{ex:pepin}
\idea If $\mathcal F_n$ is prime, Euler's criterion says $3^{(\mathcal F_n-1)/2}\equiv\left(\frac3{\mathcal F_n}\right)$, and reciprocity computes this symbol from $\mathcal F_n\bmod3$. For the converse, the congruence gives the exact order of $3$ (a power of $2$ whose half fails), and an element of order $\mathcal F_n-1$ forces $\varphi(\mathcal F_n)=\mathcal F_n-1$, i.e.\ primality. This is the template of Lucas-type primality proofs.
\writeup If $\mathcal F_n$ is prime: $\mathcal F_n\equiv1\pmod4$ and $\mathcal F_n\equiv2\pmod3$ (as $2^{2^n}\equiv1\pmod 3$, $2^n$ being even), so by reciprocity $\left(\frac3{\mathcal F_n}\right)=\left(\frac{\mathcal F_n}3\right)=\left(\frac23\right)=-1$, and Euler's criterion gives the congruence. Conversely, if $3^{(\mathcal F_n-1)/2}\equiv-1$, then $3^{\mathcal F_n-1}\equiv1$. This already forces $\gcd(3,\mathcal F_n)=1$, so $\operatorname{ord}(3)$ divides $\mathcal F_n-1=2^{2^n}$ but not its half: it equals $\mathcal F_n-1$. Since the order divides $\varphi(\mathcal F_n)\le\mathcal F_n-1$, we get $\varphi(\mathcal F_n)=\mathcal F_n-1$, so every $1\le k<\mathcal F_n$ is prime to $\mathcal F_n$, and $\mathcal F_n$ is prime.
\end{solution}

\begin{solution}{ex:gausssum}
\idea The Gauss sum $g$ is an algebraic number whose square is $\pm p$, so it ``is'' $\sqrt{p^*}$ inside $\Z[\zeta]$. Then $\left(\frac{p^*}q\right)$, which by Euler's criterion is $(p^*)^{(q-1)/2}=g^{q-1}$ modulo $q$, can be computed in a second way: raising to the power $q$ modulo $q$ is additive (the Frobenius), and it permutes the $\zeta^a$, producing $\left(\frac qp\right)$. Comparing the two computations gives reciprocity. The only subtle point is to come back from congruences in $\Z[\zeta]$ to congruences in $\Z$, which needs the linear independence of the powers of $\zeta$.
\writeup (a) With $\chi=\left(\frac\cdot p\right)$,
\[
g^2=\sum_{a,b\ne0}\chi(ab)\zeta^{a+b}.
\] Put $b=ac$ ($c$ runs over non-zero residues for fixed $a$): $\chi(ab)=\chi(a^2c)=\chi(c)$, so
\[g^2=\sum_{c\ne0}\chi(c)\sum_{a\ne0}\zeta^{a(1+c)}.\]
The inner sum is $p-1$ if $c\equiv-1$ and $-1$ otherwise (Volume~0, the roots-of-unity orthogonality proposition). Hence
\[\begin{aligned}g^2&=\chi(-1)(p-1)-\sum_{c\ne0,-1}\chi(c)\eqbr=\chi(-1)(p-1)+\chi(-1)=\chi(-1)p=p^*,\end{aligned}\]
using $\sum_{c\ne0}\chi(c)=0$.

(b) In any commutative ring, $(x+y)^q\equiv x^q+y^q$ modulo $q$ (the binomial coefficients $\binom qj$, $0<j<q$, are multiples of $q$), and by induction for sums of several terms. So $g^q\equiv\sum_a\chi(a)^q\zeta^{aq}=\sum_a\chi(a)\zeta^{aq}$, as $\chi(a)^q=\chi(a)$ ($q$ odd). Substituting $b=aq$: $\chi(a)=\chi(bq^{-1})=\chi(q)\chi(b)$ (as $\chi(q^{-1})=\chi(q)$), so the sum is $\chi(q)g$.

(c) First, $q\Z[\zeta]\cap\Z=q\Z$. The numbers $1,\zeta,\dots,\zeta^{p-2}$ are linearly independent over $\Q$: a relation would give a non-zero $R\in\Q[X]$ of degree $\le p-2$ with $R(\zeta)=0$; then $\gcd(R,\Phi_p)$, which is a $\Q[X]$-combination of $R$ and $\Phi_p$ (B\'ezout in $\Q[X]$), also vanishes at $\zeta$, so it is a non-constant proper divisor of $\Phi_p$, contradicting the irreducibility of $\Phi_p$ (Volume~0, the polynomial exercise on prime cyclotomic irreducibility). Every element of $\Z[\zeta]$ is $\sum_{i\le p-2}c_i\zeta^i$ with $c_i\in\Z$ (reduce higher powers with $\zeta^{p-1}=-1-\zeta-\dots-\zeta^{p-2}$). If an integer $n$ equals $q\sum c_i\zeta^i$, uniqueness of coordinates gives $n=qc_0$.

Now multiply (b) by $g$: $g^{q+1}\equiv\chi(q)g^2$, that is $(p^*)^{(q+1)/2}\equiv\chi(q)p^*$ modulo $q\Z[\zeta]$. Both sides are integers, so the congruence holds modulo $q$ in $\Z$. As $p^*$ is invertible modulo $q$: $(p^*)^{(q-1)/2}\equiv\chi(q)\pmod q$, and by Euler's criterion the left side is $\left(\frac{p^*}q\right)$; both are $\pm1$ and $q>2$, so they are equal. Finally $\left(\frac{p^*}q\right)=\left(\frac{-1}q\right)^{\frac{p-1}2}\left(\frac pq\right)=(-1)^{\frac{q-1}2\cdot\frac{p-1}2}\left(\frac pq\right)$, which is Theorem~\ref{thm:qr}.

(d) The number of $x$ with $x^2\equiv a$ is $1+\chi(a)$ (two roots for a residue, none for a non-residue, one for $0$), so $\sum_x\zeta^{x^2}=\sum_a(1+\chi(a))\zeta^a=\sum_a\zeta^a+g=g$. So $G^2=p^*$: $G=\pm\sqrt p$ if $p\equiv1\pmod4$ and $\pm i\sqrt p$ if $p\equiv3\pmod4$, consistent with $|G|^2=p$. A numerical calculation suggests that the sign is always $+$; this is a much harder theorem of Gauss, proved in Volume~4 with Fourier analysis.
\end{solution}

\begin{solution}{ex:t1-residue-chance}
\idea Uniform sampling converts an exact finite character sum into an
expectation. Marginal balance does not imply independence: examine the
product, rather than guessing from a picture of the signs.
\writeup Each variable is $0$ with probability $1/p$, and each of $1,-1$
with probability $(p-1)/(2p)$. Translation by $1$ permutes the sample space,
so the distributions agree. Thus $\mathbb E X=\mathbb E Y=0$ and
$\operatorname{Var}(X)=\operatorname{Var}(Y)=(p-1)/p$.
Exercise~\ref{ex:consecutive} gives
\[
\mathbb E(XY)=\frac1p\sum_{a=0}^{p-1}\chi(a(a+1))=-\frac1p.
\]
Consequently the covariance is $-1/p$. Independence would imply
$\mathbb E(XY)=(\mathbb E X)(\mathbb E Y)=0$, contradicting this value. Even a balanced
multiplicative sign pattern need not give independent neighbouring signs.
\end{solution}

\section*{Sources and further reading}

\begin{itemize}
\item K.~Ireland and M.~Rosen, \emph{A Classical Introduction to Modern Number Theory}, 2nd ed., Springer, 1990, Chapters~5--6. Reciprocity, Gauss sums and the beginning of cyclotomy.
\item F.~Lemmermeyer, \emph{Reciprocity Laws: From Euler to Eisenstein}, Springer, 2000. The full history, with a list of all known proofs of quadratic reciprocity (kept up to date on the author's web page).
\item D.~A.~Cox, \emph{Primes of the Form $x^2+ny^2$}, 3rd ed., AMS Chelsea, 2022, \S1. How reciprocity answers Fermat's questions, the theme of Chapters~7 and~10.
\item R.~Solovay and V.~Strassen, A fast Monte-Carlo test for primality, \emph{SIAM Journal on Computing} 6 (1977), 84--85. The Jacobi symbol at work in the first randomised primality test.
\item A.~Granville, It is easy to determine whether a given integer is prime, \emph{Bulletin of the American Mathematical Society} 42 (2005), 3--38. A survey of where primality testing went next, up to the AKS algorithm.
\item J.~J.~Gray, \emph{A commentary on Gauss's mathematical diary}, Expositiones Mathematicae 2 (1984), 97--130, entry of 8 April 1796; F.~Lemmermeyer, \emph{Reciprocity Laws}, Springer, 2000, whose appendix lists the published proofs.
\end{itemize}
